% Needs \usepackage{listings} and \usepackage{booktabs}.
% Source: tap_vs_csr_objdump/ (objdump/inner_loop_*.txt, results_*.txt)
\subsection{Why the Tap Table Is Faster than the CSR Table}
\label{sec:tap-vs-csr}

The tap table reaches a better inference speed than the CSR table with a much smaller table
(Table~\ref{tab:tap-vs-csr}). Both tables are placed in DTCM, which has no wait states, so the
smaller size does not make a single load faster. The gain comes from the number of loads and
instructions that are executed for each MAC. Two things reduce them:
\begin{enumerate}
    \item Strided memory accesses lower the number of metadata fetches.
    \item The compiler unrolls the output-channel loop, which removes the loop overhead.
\end{enumerate}
To check this, the two kernels were disassembled from the same builds that were used for the
measurements (GCC 14.3.1, \texttt{-O3}), with \texttt{arm-none-eabi-objdump}. The CSR kernel is
a separate function that is called for both layers. The tap kernel is inlined twice into
\texttt{SNN\_Run\_Timestep}, once for each layer.

\subsubsection*{Strided Memory Accesses}

In the CSR table, every connection stores its own output neuron and weight index. For each MAC,
the kernel first loads \texttt{out\_idx[c]} and \texttt{weight\_idx[c]}, and only then can it load
the weight and the membrane potential (Listing~\ref{lst:csr-inner}). This gives two metadata
fetches for every MAC.

\begin{lstlisting}[basicstyle=\ttfamily\scriptsize, columns=fullflexible, keepspaces=true,
    frame=single, caption={CSR kernel: index fetches and loop control, both for every MAC.},
    label={lst:csr-inner}]
// Index fetches (every MAC):
const uint16_t o = out_idx[c];
    800bf0a:  ldrh.w  r3, [r1, ip, lsl #1]
const uint16_t w = weight_idx[c];
    800bf0e:  ldrh.w  lr, [r4, ip, lsl #1]
// Weight load with the fetched index:
q31_t v = ... (q31_t)weights[w];
    800bf1a:  ldrh.w  lr, [r2, lr, lsl #1]
// Loop control (every MAC):
for (uint16_t c = row_ptr[i]; c < row_ptr[i + 1]; c++)
    800bf2a:  add.w   ip, ip, #1
    800bf32:  ldrh    r3, [r7, #2]
    800bf34:  uxth.w  ip, ip
    800bf38:  cmp     r3, ip
    800bf3a:  bhi.n   800bf0a
\end{lstlisting}

The tap table stores only the spatial part of a connection. One entry gives the output position
(\texttt{tap\_out[t]}) and the kernel position (\texttt{tap\_w[t]}), and it is fetched once.
The same entry is then used for every output channel, because the neuron and the weight of the
next output channel are at a fixed distance (stride): \texttt{out\_spatial} neurons and
\texttt{in\_ch}$\cdot$\texttt{kernel\_size} weights. In Listing~\ref{lst:tap-inner}, the two
index fetches at the top are shared by eight MACs. For the second output channel, the neuron
address is the first one plus a constant ($49 \cdot 8 = 392$\,bytes), and the weight is read with
the same index register \texttt{r2} from the next weight block. So the metadata fetches drop from
two per MAC to two per \texttt{out\_ch} MACs: $2/8$ in Conv~1 and $2/16$ in Conv~2.

\begin{lstlisting}[basicstyle=\ttfamily\scriptsize, columns=fullflexible, keepspaces=true,
    frame=single, caption={Tap kernel in Conv~1: index fetches and loop control once per tap,
    shared by the eight unrolled output channels.},
    label={lst:tap-inner}]
// Index fetches (once per tap, before the out_ch loop):
LIFNeuron* n0 = neurons + tap_out[t];
    800ccfa:  ldrh.w  r3, [r5], #2
const q15_t* w0 = w_ic + tap_w[t];
    800ccfe:  ldrb.w  r2, [r4], #1
// Neuron of each out_ch: constant stride, no index fetch:
LIFNeuron* n = &n0[(uint32_t)oc * out_spatial];
    800cd1e:  add.w   r8, r3, #392           (oc = 1)
    800cd3e:  add.w   r8, r3, #784           (oc = 2)
// Weight of each out_ch: same index r2, next weight block:
... (q31_t)w0[(uint32_t)oc * w_oc_stride];
    800cd0e:  ldrh.w  r8, [ip, r2, lsl #1]   (oc = 0)
    800cd22:  ldr     r0, [pc, #492]         (oc = 1, block address)
    800cd2a:  ldrh.w  r0, [r0, r2, lsl #1]   (oc = 1)
// Loop control (once per tap, out_ch loop unrolled):
for (uint16_t t = tap_ptr[p]; t < tap_ptr[p + 1]; t++)
    800cdfa:  ldr     r3, [r7, #56]
    800cdfc:  cmp     r3, r5
    800cdfe:  bne.w   800ccf6
\end{lstlisting}

Not all eight weight-block addresses fit in registers. Six of them are reloaded from the stack
or the literal pool (\texttt{ldr r0, [pc, \#492]} above). These are extra loads, but they read
constants and not the table. Even with them, the tap loop executes 4.1 loads per MAC in Conv~1
and 3.0 in Conv~2, compared with 6 in the CSR loop (Table~\ref{tab:tap-vs-csr}).

\subsubsection*{Loop Unrolling}

The output-channel loop of the tap kernel has a fixed length: \texttt{out\_ch} is 8 in Conv~1
and 16 in Conv~2. The kernel is inlined with these values as constants, so the compiler unrolls
the loop completely. Listing~\ref{lst:tap-inner} has no counter, compare or branch for the
output channels, and the neuron offsets ($392, 784, \dots, 2744$) are encoded directly in the
instructions. The only loop control is three instructions at the end of each tap, shared by 8
or 16 MACs.

The CSR inner loop cannot be unrolled in the same way, because its length,
\texttt{row\_ptr[i+1]}$-$\texttt{row\_ptr[i]}, is only known at run time and is different for
pixels at the edge of the input. Every MAC therefore pays five instructions of loop control
(Listing~\ref{lst:csr-inner}):
the counter increment, a zero-extension because the counter is \texttt{uint16\_t}, a re-read of
\texttt{row\_ptr[i+1]}, the compare and the branch. The bound and the input spike are read again
in every iteration because the CSR kernel is a separate function with pointer arguments, and the
compiler cannot prove that the store to the membrane potential does not change them. In the
inlined tap kernel the arrays are known, so the bound is computed once per pixel. In Conv~2 the
input spike also stays in a register, while in Conv~1 it is still re-read, because the input
comes through a pointer argument of \texttt{SNN\_Run\_Timestep}.

\subsubsection*{Effect on the Measured Cycles}

The per-sample cycles of the 694-sample runs were fitted as
$\mathrm{cycles} = a \cdot T + b \cdot \mathrm{MACs}$, where $T$ is the number of timesteps of the
sample, $a$ is the fixed cost per timestep and $b$ is the cost per MAC ($R^2 \geq 0.9999$ for
all four fits). Table~\ref{tab:tap-vs-csr} shows that the cost per MAC drops from about 13 to
about 9.2--9.5 cycles, about 28\,\% lower in both layers. This is in line with the lower number
of instructions and loads per MAC in the inner loops. The fixed cost per timestep is almost the same for both kernels, because
the decay and the soft-reset loops over all neurons are identical. For this reason the gain per
timestep is smaller than the gain per MAC: 7.9\,\% in Conv~1, where the fixed part is about
two thirds of the CSR time, and 16.8\,\% in Conv~2, where it is about one third.

\begin{table}[htbp]
\centering
\caption{CSR and tap kernels: inner loop from the disassembly, cost fitted from the 694-sample
runs, and table size.}
\label{tab:tap-vs-csr}
\begin{tabular}{lrrrr}
\toprule
 & \multicolumn{2}{c}{Conv~1} & \multicolumn{2}{c}{Conv~2} \\
\cmidrule(lr){2-3} \cmidrule(lr){4-5}
 & CSR & Tap & CSR & Tap \\
\midrule
MACs per loop iteration       & 1     & 8     & 1     & 16    \\
Instructions per MAC          & 14    & 9.5   & 14    & 8.25  \\
Loads per MAC                 & 6     & 4.1   & 6     & 3.0   \\
Metadata loads per MAC        & 2     & 0.25  & 2     & 0.125 \\
Loop-control instructions per MAC & 5 & 0.38  & 5     & 0.19  \\
\midrule
Fixed cost $a$ [cycles/timestep] & 9{,}267 & 9{,}516 & 5{,}075 & 5{,}292 \\
Cost per MAC $b$ [cycles]        & 12.98   & 9.24    & 13.16   & 9.49    \\
Mean cycles per timestep         & 13{,}929 & 12{,}835 & 14{,}677 & 12{,}217 \\
\midrule
Table size [bytes]            & 50{,}578 & 2{,}554 & 42{,}258 & 343 \\
\bottomrule
\end{tabular}
\end{table}
